\begin{table}[!htbp]
\centering
\small
\begin{threeparttable}
\caption{Population-decile gradient (relative to the largest decile)}
\label{tab:deciles}
\begin{tabular}{@{}ccc@{}}
\toprule
Population decile & Any subapplication & Any funded project \\
\midrule
1 & $-34.01$ $(2.80)$ & $-35.54$ $(2.79)$ \\
2 & $-30.75$ $(2.59)$ & $-32.36$ $(2.57)$ \\
3 & $-28.06$ $(2.54)$ & $-29.64$ $(2.53)$ \\
4 & $-25.04$ $(2.52)$ & $-26.89$ $(2.52)$ \\
5 & $-21.40$ $(2.51)$ & $-23.09$ $(2.50)$ \\
6 & $-19.92$ $(2.55)$ & $-21.13$ $(2.53)$ \\
7 & $-18.45$ $(2.41)$ & $-19.86$ $(2.40)$ \\
8 & $-14.75$ $(2.38)$ & $-15.86$ $(2.41)$ \\
9 & $-10.43$ $(2.37)$ & $-11.50$ $(2.36)$ \\
10 & 0 (reference) & 0 (reference) \\
\bottomrule
\end{tabular}
\begin{tablenotes}[flushleft]
\footnotesize
\item \textit{Notes:} Declaration fixed effects with hazard and demographic controls (Table~\ref{tab:baseline}, column 4). Coefficients multiplied by 100; county-clustered standard errors in parentheses.
\end{tablenotes}
\end{threeparttable}
\end{table}
